%% kapa: exam=1 round=37 year=2026 period=1 form=A subject=civil_law subject_ko=민법 seq=40 number=40
%% source: Q-Net 감정평가사 기출문제 (cst00309) artlSeq=5258924 file=raw/1008/5258924/2026년 제37회 감정평가사 1차 1교시.pdf
%% conversion: kordoc 4.24.51 -> tools/gen1.py (auto); review=proofread-sonnet; flags=none
\begin{question}{40}
근저당권에 관한 설명으로 옳지 않은 것은? (다툼이 있으면 판례에 따름)
\begin{choices}
  \item 후순위담보권자가 경매를 신청한 경우, 선순위 근저당권의 피담보채권은 매수인이 매각대금을 완납하여 그 근저당권이 소멸하는 때에 확정된다.
  \item 근저당권설정자가 적법하게 기본계약을 해지하면 근저당권의 피담보채권은 확정된다.
  \item 확정된 피담보채권액이 채권최고액을 초과하는 경우, 물상보증인은 채권최고액만을 변제하더라도 근저당권의 말소를 청구할 수 있다.
  \item 근저당권을 설정한 후 피담보채무가 확정되기 전에 당사자의 합의로 채무자를 변경하고자 하는 경우에는 후순위저당권자 등 이해관계인의 승낙을 받아야 한다.
  \item 후순위 근저당권자는 민법 제364조에서 규정하는 제3취득자에 해당하지 않으므로 선순위 근저당권의 피담보채무가 확정된 경우, 후순위 근저당권자는 그 확정된 피담보채무 전부를 변제하여야 비로소 선순위 근저당권의 말소를 청구할 수 있다.
\end{choices}
\end{question}
\answer{4}
